\section{Proposed Riemannian MMALS architecture}

\subsection{Architecture overview}
The proposed Riemannian Geometry for MMALS track, abbreviated \RGM{}, adds intrinsic operators at selected interfaces rather than rebuilding every host.

\begin{figure}[H]
\centering
\resizebox{\textwidth}{!}{%
\begin{tikzpicture}[
  node distance=8mm and 10mm,
  block/.style={draw,rounded corners,minimum width=26mm,minimum height=9mm,align=center,fill=white},
  geom/.style={draw,rounded corners,minimum width=29mm,minimum height=10mm,align=center,fill=lightblue,draw=mmalsblue},
  host/.style={draw,rounded corners,minimum width=21mm,minimum height=9mm,align=center,fill=lightgreen,draw=mmalsgreen},
  memory/.style={draw,rounded corners,minimum width=30mm,minimum height=10mm,align=center,fill=lightgold,draw=mmalsgold},
  arrow/.style={-{Latex[length=2mm]},thick},
  dashedarrow/.style={-{Latex[length=2mm]},thick,dashed}
]
  \node[block] (x) {Input $x$};
  \node[block,right=of x] (enc) {Controlled sensory\\encoder $z_0$};
  \node[geom,right=of enc] (context) {Context manifold\\$c=q_\eta(z_0,m)$\\RMLR optional};
  \node[geom,right=of context] (route) {Route simplex\\$r\in\Delta^{H-1}$\\Fisher-Rao};

  \node[host,below=12mm of route,xshift=-30mm] (h1) {Host $g_1$};
  \node[host,right=4mm of h1] (h2) {Host $g_2$};
  \node[host,right=4mm of h2] (h3) {Host $g_H$};
  \node[block,below=12mm of h2] (synth) {Fungal transforms $T_i$\\and synthesis $z_f$};
  \node[geom,right=of synth] (spd) {Functional descriptor\\$C$ or $R$\\SPD/correlation};
  \node[block,right=of spd] (head) {Prediction / control\\$\hat y$};

  \node[memory,below=14mm of spd] (mem) {Dual memory\\samples + geometric state\\$C^\star,r^\star$};
  \node[geom,left=of mem] (metric) {Metric controller\\fixed LE / Cholesky /\\learnable pullback};

  \draw[arrow] (x)--(enc);
  \draw[arrow] (enc)--(context);
  \draw[arrow] (context)--(route);
  \draw[arrow] (route.south) -- ++(0,-4mm) -| (h1.north);
  \draw[arrow] (route.south) -- ++(0,-4mm) -| (h2.north);
  \draw[arrow] (route.south) -- ++(0,-4mm) -| (h3.north);
  \draw[arrow] (h1)--(synth);
  \draw[arrow] (h2)--(synth);
  \draw[arrow] (h3)--(synth);
  \draw[arrow] (synth)--(spd);
  \draw[arrow] (spd)--(head);
  \draw[dashedarrow] (mem)--(spd);
  \draw[dashedarrow] (mem.west)--(metric.east);
  \draw[dashedarrow] (metric.north)--(synth.south);
  \draw[dashedarrow] (mem.north) to[bend left=15] (route.south east);

  \node[draw=mmalsred,dashed,rounded corners,fit=(context)(route)(h1)(h3)(synth)(spd)(metric)(mem),inner sep=5mm,label={[text=mmalsred]below:Geometry is admitted component by component, not assumed globally}] {};
\end{tikzpicture}%
}
\caption{Proposed selective Riemannian extension of \MMALS{}. The sensory encoder and hosts may remain ordinary neural networks. Intrinsic geometry is introduced at the context, route, functional-descriptor, and memory-transport interfaces only when ablations justify it.}
\label{fig:architecture}
\end{figure}

\subsection{Context inference with intrinsic classification}
The mature model infers context without a task identifier. Let $c\in\cM_c$ be a structured context representation and $\mu_k\in\cM_c$ context prototypes. A Euclidean classifier applies a linear map to coordinates. Riemannian MLR instead constructs logits from intrinsic geometry, requiring a logarithmic map and a tangent normal $A_k\in T_{\mu_k}\cM_c$ \citep{chen2024rmlr,chen2026thesis}. A generic score can be written schematically as
\begin{equation}
  \ell_k(c)
  = b_k + \langle A_k,\Log_{\mu_k}(c)\rangle_{\mu_k}
  \cdot \omega_k(c),
  \label{eq:rmlr-context}
\end{equation}
where $\omega_k$ encodes the Riemannian-trigonometric normalization required by the selected formulation. The operational rule is strict: RMLR is tested only after the context representation has a declared manifold and valid logarithmic map. It is not used as decorative terminology for an unconstrained vector.

\subsection{Intrinsic normalization at stable interfaces}
Sequential training can shift the mean and variance of structured states. LieBN and GyroBN offer principled normalization when the state manifold admits the corresponding algebraic structure \citep{chen2024liebn,chen2025gyrobn}. In \RGM{}, candidate insertion points are:
\begin{enumerate}[leftmargin=1.7em]
  \item after an SPD/correlation pooling layer used to build $C_k$ or $R_k$;
  \item before intrinsic classification of context or functional state;
  \item inside a host only when that host natively produces manifold-valued features.
\end{enumerate}
Normalization is not placed on the route simplex by analogy alone. The route has its own statistical geometry and conservation constraints.

\subsection{Hyperbolic geometry only for demonstrated hierarchy}
Hyperbolic spaces are useful for hierarchical or tree-like relations because negative curvature provides rapidly growing volume. Chen's Proper Velocity model offers an unconstrained hyperbolic representation and complete neural layers intended to mitigate instability near constrained-model boundaries \citep{chen2026pvnn}. For \MMALS{}, a hyperbolic context or host space is admitted only if:
\begin{enumerate}[leftmargin=1.7em]
  \item the benchmark contains a grounded hierarchy or branching regime structure;
  \item hyperbolic distortion beats Euclidean and spherical controls;
  \item the hierarchy improves routing, transfer, or sample efficiency rather than only visualization;
  \item numerical stability and compute are recorded.
\end{enumerate}
RotatedMNIST has a circular factor and is therefore a poor justification for hyperbolic geometry. A hierarchical engineering-evidence graph, taxonomy, or compositional task family is a better later benchmark.


\section{Memory transport and objective}

\subsection{Transporting functional memory}
A stored tangent direction $V_t\in T_{C_t}\cM_f$ cannot be directly compared with a direction at $C_{t+1}$ because the tangent spaces differ. Parallel transport provides a principled alignment:
\begin{equation}
  \tilde V_{t\to t+1}
  = \PT_{C_t\to C_{t+1}}(V_t).
  \label{eq:parallel-transport}
\end{equation}
A transport consistency loss is
\begin{equation}
  \cL_{\mathrm{transport}}
  = \E_k\left[
    \left\|V_{k,t+1}-\PT_{C_{k,t}\to C_{k,t+1}}(V_{k,t})\right\|_{C_{k,t+1}}^2
  \right].
  \label{eq:transport-loss}
\end{equation}
This must be distinguished from merely penalizing movement. Healthy adaptation can move along the manifold while preserving functional relations. The test is whether transported directions remain predictive of behavior.

\subsection{Complete training objective}
A selective geometry-aware objective is
\begin{align}
  \cL_{\RGM}
  =&\ \cL_{\mathrm{task}}
  + \lambda_y \cL_{\mathrm{distill}}
  + \lambda_r \cL_{r,\mathrm{FR}}
  + \lambda_f d_\eta(C,C^\star)^2 \\
  &+ \lambda_T \cL_{\mathrm{transport}}
  + \lambda_n \cL_{\mathrm{norm}}
  + \lambda_m \cL_{\mathrm{metric}}
  + \lambda_c \cC_{\mathrm{carbon}}
  + \lambda_g \cL_{\mathrm{gain}}.
  \label{eq:full-objective}
\end{align}
The terms mean:
\begin{itemize}[leftmargin=1.5em]
  \item $\cL_{\mathrm{task}}$: current supervised, self-supervised, or control objective;
  \item $\cL_{\mathrm{distill}}$: output behavior preservation;
  \item $\cL_{r,\mathrm{FR}}$: intrinsic route memory;
  \item $d_\eta(C,C^\star)^2$: functional-state preservation under fixed or learned geometry;
  \item $\cL_{\mathrm{transport}}$: consistency of functional directions through time;
  \item $\cL_{\mathrm{norm}}$: optional intrinsic normalization statistics;
  \item $\cL_{\mathrm{metric}}$: complexity, conditioning, and deviation control for the learned metric;
  \item $\cC_{\mathrm{carbon}}$: active-route and computation cost;
  \item $\cL_{\mathrm{gain}}$: mutualistic contribution constraint estimated with held-out or rolling evidence.
\end{itemize}

Not all terms are activated simultaneously. The ablation ladder starts from output distillation and introduces exactly one geometric mechanism at a time.

\subsection{Two basic propositions}

\begin{proposition}[Route preservation does not imply functional preservation]
There exist teacher and current models with identical route distributions and different synthesized functions.
\end{proposition}
\begin{proof}
Take $r^\star=r=(1,0,\ldots,0)$ and choose $T_1^\star(g_1^\star(x))\neq T_1(g_1(x))$. Then the route distance is zero under any valid metric on the simplex, while $z_f^\star\neq z_f$. Therefore route memory is not sufficient for functional memory.
\end{proof}

\begin{proposition}[Pullback distance is Euclidean after a valid chart]
Let $\phi:\cM\rightarrow\R^d$ be a diffeomorphism and $g=\phi^*g_E$. Then
\begin{equation}
  d_g(p,q)=\|\phi(p)-\phi(q)\|_2
\end{equation}
for the pullback Euclidean geometry.
\end{proposition}
\begin{proof}
The diffeomorphism is an isometry by construction of the pullback metric. Curve lengths and the minimizing geodesic distance are therefore preserved under $\phi$.
\end{proof}

\begin{remark}
The proposition does not mean that any learned embedding is a valid Riemannian chart. Smoothness, invertibility, numerical conditioning, and domain coverage are implementation obligations. A black-box encoder with no inverse does not automatically define the tractable pullback geometry used here.
\end{remark}

